% Auto-generated by python -m src.q2_learning_curve --model masked_lr
% Upload q22_learning_curve.png to Overleaf alongside the report source.

% Q2.1
\begin{tabularx}{\linewidth}{@{}rrrr@{}}\toprule
Training rows & Training AUROC & Validation AUROC & Notes \\ \midrule
500 & 1.0000 & 0.9869 & Nested balanced subset, seed 42 \\
1,000 & 1.0000 & 0.9903 & Nested balanced subset, seed 42 \\
2,000 & 1.0000 & 0.9926 & Nested balanced subset, seed 42 \\
5,000 & 0.9998 & 0.9944 & Nested balanced subset, seed 42 \\
8,000 & 0.9998 & 0.9949 & Nested balanced subset, seed 42 \\
\bottomrule\end{tabularx}
\smallnote{Configuration held fixed at masked word-pattern 1--3-gram TF--IDF (up to 200{,}000
features) with logistic regression, $C=4.0$. Subsets are nested, class-balanced prefixes drawn from
the shared 8{,}000-row training partition; all points use the same 2{,}000-row validation set.}

% Q2.2
\begin{figure}[ht]
  \centering
  \includegraphics[width=0.78\linewidth]{q22_learning_curve.png}
  \caption{Training and validation AUROC versus labeled training-set size. All points use the same
  fixed 2,000-row validation set and identical hyperparameters.}
  \label{fig:q2-learning-curve}
\end{figure}

% Q2.3
Using one fixed masked n-gram logistic-regression configuration, training AUROC is 1.0000 with 500 examples and 0.9998 with 8,000, while validation AUROC changes from 0.9869 to 0.9949. The training-validation gap therefore changes from 0.0131 to 0.0049. High training performance together with a lower validation score at small sample sizes indicates overfitting rather than underfitting: the model can fit the available rows, but its estimates generalize less reliably. The largest validation gain occurs between 500 and 1,000 examples. The final increment is small, indicating diminishing returns from additional labels under this representation. These conclusions apply to the fixed in-distribution split and may not fully predict performance on shifted Kaggle test data.
